\subsection{Language Modeling}\label{sec:language-modeling}

Before studying \textbf{how embeddings are learned}, we need to understand the \hl{problem that motivated their development}. Historically, word embeddings did not emerge as an isolated idea. They were discovered while researchers were trying to solve a much broader problem: \textbf{\emph{Language Modeling}}. That is: \emph{given a sequence of words, can we predict how likely that sequence is, and what word should come next?}

\longline

\subsubsection{The Language Modeling Problem}\label{sec:the-language-modeling-problem}

First of all, a \definition{Language Model (LM)} is a \textbf{probabilistic model that assigns probabilities to sequences of words}. Given a sentence:
\begin{equation*}
    s = \left(w_1, w_2, \dots, w_k\right)
\end{equation*}
We want to estimate the probability of that sentence:
\begin{equation*}
    P(s) = P\left(w_1, w_2, \dots, w_k\right)
\end{equation*}
Where $w_i$ is the $i$-th word in the sentence. This quantity measures how likely the sentence is according to the language.

\highspace
\begin{flushleft}
    \textcolor{Green3}{\faIcon{question-circle} \textbf{Why do we need Language Models?}}
\end{flushleft}
Many Natural Language Processing (NLP) applications depend on language modeling: Chatbots, Machine Translation, Speech Recognition, Information Retrieval, Document Classification, Text Generation, Large Language Models (LLMs). The performance of many real-world systems depends on how language is represented and modeled.

\highspace
For example, consider two sentences:
\begin{enumerate}
    \item \emph{The cat is sleeping on the sofa.}
    \item \emph{Sofa sleeping cat the on is.}
\end{enumerate}
A human immediately recognizes that:
\begin{equation*}
    P(\text{The cat is sleeping on the sofa}) \gg P(\text{Sofa sleeping cat the on is})
\end{equation*}
The first sentence is a valid English sentence, while the second one is not. A language model \textbf{attempts to learn exactly this distinction}. But now the question is: \emph{\textbf{how can we build a model that learns to assign high probabilities to valid sentences and low probabilities to invalid ones?}}

\highspace
\begin{flushleft}
    \textcolor{Green3}{\faIcon{link} \textbf{Probability of Sentences}}
\end{flushleft}
At first glance, estimating the probability of a sentence:
\begin{equation*}
    P(s_k) = P\left(w_1, w_2, \dots, w_k\right)
\end{equation*}
\textbf{Seems impossible}. Even a short sentence can contain many words, and the number of possible sentences is astronomical. Fortunately, probability theory gives us a useful tool: the \definition{Chain Rule of Probability}. \hl{To compute the probability of several events happening together, we can decompose it into a product of conditional probabilities.} For example, for two events $A$ and $B$:
\begin{equation*}
    P\left(A,B\right) = P\left(A\right) P\left(B \, | \, A\right)
\end{equation*}
For three events $A$, $B$, and $C$:
\begin{equation*}
    P\left(A,B,C\right) = P\left(A\right) P\left(B \, | \, A\right) P\left(C \, | \, A, B\right)
\end{equation*}
In general, for $n$ events $A_1, A_2, \dots, A_n$:
\begin{equation}\label{eq:chain-rule-of-probability}
    \begin{aligned}
        P\left(A_1, A_2, \dots, A_n\right) &= P\left(A_1\right) P\left(A_2 \, | \, A_1\right) \cdots P\left(A_n \, | \, A_1, A_2, \dots, A_{n-1}\right) \\
        &= \prod_{i=1}^n P\left(A_i \, | \, A_1, A_2, \dots, A_{i-1}\right)
    \end{aligned}
\end{equation}
It is called the \textbf{\emph{Chain Rule of Probability}} because it allows us to break down a complex joint probability into simpler conditional probabilities. Applying the chain rule to language modeling, we can express the probability of a sentence as:
\begin{equation*}
    P\left(w_1, w_2, \dots, w_k\right) = P\left(s_k\right) = P\left(w_1\right) P\left(w_2 \, | \, w_1\right) \cdots P\left(w_k \, | \, w_1, w_2, \dots, w_{k-1}\right)
\end{equation*}
This means that to compute the probability of a sentence, we can \textbf{compute the probability of each word given the previous words}. For example, the sentence ``\emph{the cat sleeps}'' can be decomposed as:
\begin{equation*}
    P\left(\text{the cat sleeps}\right) = P\left(\text{the}\right) P\left(\text{cat} \, | \, \text{the}\right) P\left(\text{sleep} \, | \, \text{the cat}\right)
\end{equation*}

\begin{flushleft}
    \textcolor{Green3}{\faIcon{question-circle} \textbf{With the Chain Rule, have we actually solved anything?}}
\end{flushleft}
With the Chain Rule, we have reduced the problem of estimating the probability of a sentence to estimating the probability of each word given the previous words. However, \emph{have we actually solved the problem?} After all, we \textbf{still need to estimate probabilities}. The answer is that the Chain Rule does \textbf{not solve the language modeling problem}. It only \hl{reformulates it into a more structured form.}

\highspace
Instead of asking ``\emph{what is the probability of every possible sentence?}'' we ask ``\emph{what is the probability of a word appearing after a given context?}''. For example:
\begin{equation*}
    P\left(\text{sleep} \, | \, \text{the cat}\right)
\end{equation*}
Can be interpreted as: \emph{given the context ``the cat'', how likely is the next word ``sleeps''?} This transforms language modeling into a collection of \textbf{next-word prediction problems}.

\highspace
The chain rule reveals that \textbf{if we can estimate probabilities} of the form:
\begin{equation*}
    P\left(w_i \, | \, w_1, w_2, \dots, w_{i-1}\right)
\end{equation*}
Then we can \hl{automatically compute the probability of any sentence.} Consequently, researchers shifted their attention from sentence probabilities to next-word prediction. For example, given the context:
\begin{center}
    \emph{The cat is}
\end{center}
A language model might estimate the following probabilities for the next word:
\begin{align*}
    P\left(\text{sleeping} \, | \, \text{The cat is}\right) &= 0.70 \\
    P\left(\text{eating} \, | \, \text{The cat is}\right) &= 0.20 \\
    P\left(\text{flying} \, | \, \text{The cat is}\right) &= 0.001
\end{align*}
The most likely continuation is therefore: \emph{sleeping}. Notice that the objective is not really to generate text. The \textbf{objective is to estimate conditional probabilities}. \hl{Text generation is simply a consequence}: once the model knows the probabilities of possible next words, it can choose one and continue the sentence.

\highspace
But again, we have not yet solved the language modeling problem. We have only transformed it into a more structured task: estimating probabilities of the form:
\begin{equation*}
    P\left(w_i \, | \, w_1, w_2, \dots, w_{i-1}\right)
\end{equation*}
The natural question is now: \emph{\textbf{how can we estimate these probabilities in practice?}}

\highspace
\begin{flushleft}
    \textcolor{Green3}{\faIcon{link} \textbf{Connection with LLMs}}
\end{flushleft}
The problem introduced here (i.e., estimating the probability of a word given its context) is essentially the same problem solved by modern Large Language Models (LLMs). When ChatGPT generates text, it is repeatedly estimating:
\begin{equation*}
    P\left(w_i \, | \, w_1, w_2, \dots, w_{i-1}\right)
\end{equation*}
And choosing the next token according to that distribution. The difference is that:
\begin{itemize}
    \item Classical language models use simple statistical methods.
    \item Neural Language Models (our focus) use neural networks to estimate probabilities.
    \item Modern LLMs use transformers with billions of parameters.
\end{itemize}
But the underlying objective remains the same: \textbf{predict the next word (or token) given the previous context}.

\highspace
\begin{flushleft}
    \textcolor{Red2}{\faIcon{exclamation-triangle} \textbf{Why is this difficult?}}
\end{flushleft}
At this point we might think: \emph{\textbf{great, we just need to estimate conditional probabilities}} (i.e., the probability of a word given its context). Unfortunately, there is a \textbf{problem}. To compute:
\begin{equation*}
    P\left(w_i \, | \, w_1, w_2, \dots, w_{i-1}\right)
\end{equation*}
We may need to consider the entire history of the sentence. In general, for a sentence of length $k$, the context can become arbitrarily large, because the information needed to predict the next word may depend on words that appeared much earlier in the sentence. This \hl{challenge led researchers to develop $N$-gram Language Models}, which \hl{approximate the problem by considering only a limited number of previous words.}